% Supplementary Information for "Time-domain noise as an evaluation axis for raw-signal nanopore
% mapping". Typeset from supplementary.md, which is the canonical text; the two must carry the
% same numbers in the same order, which scripts/verify_dwell_journal_numbers.py checks.
%
% S1-S3 hold the values that no captured run backs. Each is marked \prov, and the script
% allows such values here and nowhere in the main manuscript. S4-S5 hold backed results that the
% main text has no room to tabulate; the script rebuilds every row of Tables S1-S3 from data.
%
% Build: latexmk -pdf supplementary.tex

\documentclass[11pt,a4paper]{article}

\usepackage[T1]{fontenc}
\usepackage[utf8]{inputenc}
\usepackage{newtxtext,newtxmath}
\usepackage[a4paper,margin=2.5cm]{geometry}
\usepackage{microtype}
\usepackage{authblk}
\usepackage{booktabs}
\usepackage[font=small,labelfont=bf,labelsep=period]{caption}
\usepackage[hidelinks]{hyperref}

\setlength{\parskip}{0pt}
\renewcommand{\Affilfont}{\small}
\renewcommand{\thesection}{S\arabic{section}}
\renewcommand{\thetable}{S\arabic{table}}
\newcommand{\prov}{\textsuperscript{\dag}}
% Repository paths in typewriter that may break at / and _, which \texttt cannot.
\DeclareUrlCommand\file{\urlstyle{tt}}

\title{\textbf{Supplementary Information}\\[6pt]
  \large Time-domain noise as an evaluation axis for raw-signal nanopore mapping}
\author[1,*]{Khushal Singh}
\author[1]{Ritu Garg}
\affil[1]{Department of Computer Engineering, National Institute of Technology Kurukshetra,
  Kurukshetra, Haryana 136119, India}
\affil[*]{To whom correspondence should be addressed: \texttt{khushal6002@gmail.com}}
\date{}

\begin{document}

\maketitle

\noindent This supplement has three parts. Sections S1--S3 report three observations made during
this work that rest in part on a measurement for which no run output was captured. Those values are
marked with a dagger (\dag). None of them is used in any claim in the main text, and
\file{scripts/verify_dwell_journal_numbers.py} lists each as outstanding, together with the run that
would settle it. Sections S4 and S5 give committed results in more detail than the main text does,
and every value in them is recomputed by the same script from files in the repository. Section S6
documents the use of AI tools in this work. Section and table numbers without an S refer to the
main text.

\section{Live accept rates against Icarust}\label{sec:s-icarust}

With the same client, reference and run length, the fixed-width arm accepted 43.5\%\prov{} of reads
and the detected arm 9.6\%\prov{} against stock Icarust, and 0.5\%\prov{} and 9.6\%\prov{} against
the patched simulator at CV 0.35. Taken at face value, stock Icarust reports the fixed-width arm as
4.5-fold\prov{} better, and the patched simulator reverses the order, which is the reversal of
Section~3.1 observed on the live interface. These are accept rates on the live loop and are not
comparable with the offline true-positive rates (TPR) of Table~2. Re-running stock and patched
Icarust against the client would settle them.

\section{Calibration target for event detection}\label{sec:s-calibration}

The calibrated detector, threshold 3.0, emits 1562 events per on-target read at CV 0.44, against an
ideal of approximately 2000, one event per base, which is 21.9\% under-segmentation; seed 42 gives
14.2\% TPR in the captured log. In a development run that was not captured, threshold 2.0 gave
2017\prov{} events per read, close to the ideal, and 2.0\%\prov{} TPR, 7.1-fold\prov{} lower. An
argument offered during development attributed the difference to frame displacement: a spurious
boundary shifts every later key in the read. That argument does not distinguish the two error
types, because a missed boundary removes an event and shifts every later event index by one exactly
as a spurious boundary adds one. A sufficient explanation is simpler. The event count measures net
error, insertions minus deletions, and is blind to errors that cancel, so a detector that emits the
ideal number of events may still place many boundaries wrongly. The recommendation that follows
does not depend on the uncaptured values: calibrate a detector against agreement or indel fraction
(Section~2.6), not against its event count.

\section{Calibration transfer between generators}\label{sec:s-transfer}

A detector threshold calibrated against signal from a generator written for this project, applied
unchanged to Squigulator signal, gave 0.6\%\prov{} TPR at CV~0, where the Squigulator-calibrated
detector gives 18.4\% [17.0--19.8] (Table~2), a 30.7-fold\prov{} difference. The two generators
differ in their amplitude-noise models, and the detector threshold is a t-statistic whose operating
point depends on that model. Whatever the magnitude, a result obtained with event detection is
interpretable only together with a statement of how the detector was calibrated and on what signal.

\section{RawHash2 at matched operating points}\label{sec:s-matched}

RawHash2's results in Tables~3 and~5 are at its default operating point, which maps some control
reads. Tables~\ref{tab:s-matched} and~\ref{tab:s-matched-cost} place it at two operating points set
by a false-positive budget, from the output of the dwell $\times$ window grid (Section~2.5) and
without new runs. \file{scripts/rawhash_zero_fp.py} reads that output from
\file{results/rawhash-grid-raw/raw.tar.gz}. It first scores every cell with no threshold, by the
rule of Section~2.5, and requires the result to reproduce \file{results/rawhash-grid-seeds.csv}
exactly, which all 54 cells do. It then thresholds the mapping quality that RawHash2 reports for
each mapping, the twelfth column of its PAF output. At zero false positives (Zero FP) the threshold
is one above the highest mapping quality given to any control read. With a 1\% budget (1\% FP) it is
the lowest threshold that leaves at most 5 of the 500 control reads mapped, so it is set by the
sixth-highest-scoring control read. An on-target read counts if it is mapped correctly with a
mapping quality at or above the threshold. Each seed's thresholds and rates are in
\file{results/rawhash-grid-zero-fp.csv}, which the verification script recomputes from the raw
output on every run.

\begin{table}[htbp]
\centering
\caption{RawHash2's correct-mapping rate (\%) at matched operating points, by dwell CV and minimizer
window $w$: the percentage of on-target reads mapped correctly at or above the lowest
mapping-quality threshold that admits no control read (Zero FP) or at most 1\% of them (1\% FP).
Mean [range] over three seeds, 500 on-target and 500 control reads per cell.}
\label{tab:s-matched}
\begin{tabular}{llccc}
\toprule
Operating point & Dwell CV & $w = 0$ & $w = 5$ & $w = 10$ \\
\midrule
Zero FP & 0.00 & 86.3 [85.2--87.0] & 79.7 [76.8--84.8] & 70.6 [67.4--73.4] \\
Zero FP & 0.10 & 88.7 [86.4--91.2] & 82.8 [77.8--85.6] & 70.1 [63.2--75.2] \\
Zero FP & 0.20 & 89.5 [87.4--92.8] & 80.3 [76.6--82.6] & 67.7 [64.2--73.6] \\
Zero FP & 0.30 & 83.1 [81.6--85.8] & 77.8 [75.2--82.6] & 65.3 [63.0--68.4] \\
Zero FP & 0.44 & 80.5 [76.2--84.6] & 78.1 [75.0--82.8] & 62.5 [60.2--66.8] \\
Zero FP & 0.60 & 82.1 [79.4--86.4] & 71.7 [70.4--73.8] & 52.5 [41.8--62.8] \\
\midrule
1\% FP & 0.00 & 90.6 [89.8--91.6] & 86.8 [85.6--88.2] & 77.2 [75.8--78.8] \\
1\% FP & 0.10 & 91.4 [90.6--92.8] & 86.4 [84.6--87.6] & 78.5 [78.2--78.6] \\
1\% FP & 0.20 & 91.1 [90.0--93.0] & 86.0 [85.4--86.8] & 76.7 [73.6--78.4] \\
1\% FP & 0.30 & 89.5 [87.0--92.6] & 84.7 [83.8--86.0] & 74.2 [70.8--76.6] \\
1\% FP & 0.44 & 88.4 [87.6--89.0] & 82.5 [81.2--84.2] & 71.9 [70.2--73.6] \\
1\% FP & 0.60 & 87.5 [86.2--89.0] & 77.9 [76.6--79.4] & 63.0 [61.2--65.8] \\
\bottomrule
\end{tabular}
\end{table}

\begin{table}[htbp]
\centering
\caption{Accuracy cost of minimizer sketching in RawHash2 at the operating points of
Table~\ref{tab:s-matched}: the rate at $w = 0$ minus that at the stated window, in percentage
points, taken within each seed; mean [range] over three seeds. A negative value means that the
sketched index mapped more reads correctly in that seed.}
\label{tab:s-matched-cost}
\begin{tabular}{llcc}
\toprule
Operating point & Dwell CV & $w = 5$ & $w = 10$ \\
\midrule
Zero FP & 0.00 & 6.7 [0.4--10.0] & 15.7 [13.4--19.6] \\
Zero FP & 0.10 & 5.9 [2.8--8.6] & 18.6 [13.2--23.2] \\
Zero FP & 0.20 & 9.1 [4.8--11.6] & 21.8 [13.8--27.6] \\
Zero FP & 0.30 & 5.3 [3.2--6.4] & 17.9 [13.6--21.4] \\
Zero FP & 0.44 & 2.5 [$-$6.6--8.2] & 18.0 [9.4--24.0] \\
Zero FP & 0.60 & 10.4 [9.0--12.6] & 29.5 [16.6--38.6] \\
\midrule
1\% FP & 0.00 & 3.8 [3.4--4.2] & 13.4 [12.8--14.6] \\
1\% FP & 0.10 & 5.0 [3.2--6.0] & 12.9 [12.0--14.6] \\
1\% FP & 0.20 & 5.1 [3.6--7.2] & 14.4 [11.8--19.4] \\
1\% FP & 0.30 & 4.7 [2.6--8.8] & 15.3 [10.4--21.8] \\
1\% FP & 0.44 & 5.9 [3.4--7.4] & 16.5 [15.0--17.4] \\
1\% FP & 0.60 & 9.6 [8.6--10.6] & 24.5 [21.4--27.0] \\
\bottomrule
\end{tabular}
\end{table}

Three features of these tables bear on the main text. First, a zero-false-positive threshold rests
on the single highest-scoring control read, and the rates at it are correspondingly noisy: at
$w = 0$ the Zero FP rate is not monotone in CV, and at CV 0.44 one seed's $w = 5$ index maps
6.6 points more of its reads correctly than its $w = 0$ index, which no seed does with the 1\% budget
or at the default (Table~\ref{tab:s-perseed}). Second, at $w = 10$ sketching costs more at both
matched operating points than at the default, at every CV (Tables~5 and~\ref{tab:s-matched-cost}).
Third, the dependence of that cost on dispersion (Section~3.3) holds with the 1\% budget at both
windows, in every seed and with disjoint ranges between CV~0 and CV 0.60; at zero false positives it
holds in every seed at $w = 10$ only.

\section{Per-seed sketching cost at the default operating point}\label{sec:s-perseed}

Table~\ref{tab:s-perseed} gives each seed's cost of sketching at RawHash2's default operating
point, the values that Table~5 summarises and Figure~4 draws: the correct-mapping rate at $w = 0$
minus that at the stated window, from \file{results/rawhash-grid-seeds.csv}. Every window of one
seed maps the same reads, so each entry is a paired difference. In each seed the cost at CV 0.60
exceeds the cost at CV~0 at both windows, which is the claim of Section~3.3. Only seed 44's cost
rises at every step of the sweep, at both windows, which is why that section does not claim a rise
at every step.

\begin{table}[htbp]
\centering
\caption{Accuracy cost of minimizer sketching in RawHash2 at its default operating point, in each
seed: correct-mapping rate at $w = 0$ minus that at the stated window, in percentage points.
Table~5 gives the mean and range of each set of three.}
\label{tab:s-perseed}
\begin{tabular}{lcccccc}
\toprule
& \multicolumn{3}{c}{$w = 5$} & \multicolumn{3}{c}{$w = 10$} \\
\cmidrule(lr){2-4}\cmidrule(lr){5-7}
Dwell CV & Seed 42 & Seed 43 & Seed 44 & Seed 42 & Seed 43 & Seed 44 \\
\midrule
0.00 & 5.0 & 5.0 & 4.0 & 12.8 & 11.8 & 10.4 \\
0.10 & 6.4 & 5.2 & 4.6 & 13.0 & 11.4 & 10.6 \\
0.20 & 5.2 & 3.6 & 5.2 & 12.4 & 10.4 & 10.8 \\
0.30 & 6.2 & 4.0 & 5.8 & 15.0 & 10.2 & 11.6 \\
0.44 & 7.2 & 5.4 & 6.6 & 15.2 & 16.4 & 14.6 \\
0.60 & 10.4 & 7.2 & 8.6 & 21.0 & 18.2 & 19.2 \\
\bottomrule
\end{tabular}
\end{table}

\section{Use of AI tools}\label{sec:s-ai}

The authors used Claude (Anthropic) through Claude Code, a tool that reads and edits files and runs
commands in the authors' environment at their direction; the models were Claude Opus 5 and Claude
Opus 5.5. The repository's history records the extent of the use: every commit in it carries a
Claude co-author line, including those that add the pipeline's source code, the analysis scripts and
the dissertation drafts. Its tasks for this article were the following.

\paragraph{Code and analysis.} It helped write the pipeline's source code and the scripts that map
and score reads, sweep dwell dispersion and draw the figures. While this article was prepared it ran
the three-seed RawHash2 grid of Section~2.5 (\file{scripts/rerun_rawhash_grid.sh}) on simulated
reads generated earlier in the work, and it wrote and ran \file{scripts/rawhash_zero_fp.py}, which
computes Section~\ref{sec:s-matched} from that grid's output. No reads were simulated while this
article was prepared.

\paragraph{Verification.} It wrote \file{scripts/verify_dwell_journal_numbers.py}, which recomputes
every number in the article and in this supplement from the committed result files, and tested the
script by planting errors in copies of both documents, each of which the script had to reject.
Checking the dissertation chapter's numbers against the raw data in this way led to several
corrections, including the replication over three seeds of a result first reported from one.

\paragraph{References.} It identified candidate references and checked each against Crossref, the
publisher's page or arXiv.

\paragraph{Text and figures.} It drafted this article from Chapter~6 of the first author's
dissertation, which it had also helped draft, and typeset the LaTeX source. Figure~1 is a schematic
drawn by a script it wrote.

The authors directed each task, reviewed the code, the data and the text, and take full
responsibility for the content. No AI tool is an author.

\end{document}
